\scanpage{48}

6.1.1966

Dear Coates,

Here a\inserted{\ }few more comments to my talk on the conjectures. The following proposition shows that the conjecture $C_{\unsure{\ell|\chi}}(X)$ is independant of the choosen polarization, and has also some extra interest, in showing the part playd by the fact that $H^i(X)$ should be ``motive-theoretically'' isomorphic to its natural dual $H^{2n-i}(X)$ (as usual, I drop the twist for simplicity).

\uline{Proposition}\quad The condition $C_{\unsure{\ell|1}}(X)$ is equivalent also to \inserted{each of} the following conditions:

a) $D_{\unsure{\ell|\chi}}(X)$ horlds, and for every $i\inserted{{}<n}$, there exists an isomorphism $H^{2n-i}(X)\to H^i(X)$ which is algebraic (i.e. induced by an algebraic correspondance class; we do not make any assertion on what it induces in degrees different from $2n-i$).

b) For every endomorphism $H^i(X)\mathrel{-}H^i(X)$ which is algebraic, the \unsure{coefficients} of the characteristic polynomial are rational, and for every $i\inserted{{}<n}$ there exists an isomorphism $H^{2n-i}(X)\mathrel{-}H^i(X)$ which is algebraic.

Proof. I sketched already how $D_{\unsure{\ell|\chi}}(X)$ implies the fact that for an algebraic endomorphism of $H^i(X)$, the coefficients of the characteristic polynomial are rational numbers, Therefore we know that a) implies b), and of course $C_{\unsure{\ell|\chi}}(X)$ implies a). It remains to prove that b) implies $C_{\unsure{\ell|\chi}}(X)$. Let $u:H^{2n-i}(X)\mathrel{-}H^i(X)$ be the given isomorphism which is algebraic, and $v:H^i(X)\mathrel{-}H^{\unsure{2n-i}}(X)$ \unsure{the} algebraic isomorphism in the opposite direction, induced by $L_X^{n-i}$. Then $uv=w$ is an automorphism of $H^i(X)$ which is algebraic, and the Hamilton-Cayley formula
\[
\unsure{u|w}^{\unsure{h|b}}-\sigma_1(w)u^{\unsure{h|b}-1}+\ldots+\deleted{(-1)}(-1)^b\sigma_b(w)=0
\]
(where the $\sigma_i(w)$ are the coefficients of the car. pol. of $w$) whos that $w^{-1}$ is a linear combination of the $w^i$, with coefficients of the type $\pm\sigma_i(w)/\sigma_b(w)$\quad (NB $b=\operatorname{rank}H^i$). The assumption implies

\mspage{47}
